    % ==== Natural Language Processing and Voice Assistants =============================================================================
    \chapter{Natural Language Processing and Voice Assistants}
    \section{Natural Language Processing}
    \subsection{Introduction to NLP}
    \subsubsection{What is NLP?}
    \lipsum[25][1-3]
    \subsubsection{Syntax, Semantics and Prosodics}
    \lipsum[26][1-3]
    \subsubsection{Phonetics and Speech}
    \lipsum[1][1]
    \subsubsection{Evaluation of NLP Systems}
    \lipsum[1][1]
    \subsection{Text Processing}
    \subsubsection{Word Vectors and Word Embeddings}
    \subsubsection{Regular Expressions}
    \subsubsection{Statistical Approaches}
    \subsubsection{Recurrent Neural Network based Approaches}
    \subsubsection{Transformer based Approaches}
    \subsection{Speech Processing}
    \subsubsection{Statistical Speech Recognition and Synthesis}
    \subsubsection{Speech Recognition and Synthesis with Deep Learning}
    \subsection{Application Scenarios}
    \subsubsection{Speech Recognition, Speech Synthesis and Machine Translation}
    \subsubsection{Information Extraction and Text Understanding}
    \subsubsection{Chatbots and Voice Assistants}
    \subsubsection{NLP in Education}
    \lipsum[26][1-3]

    \subsubsection{NLP with Python}
    \subsection{Challenges in NLP}
    \subsubsection{Data for NLP}
    \subsubsection{Domain and Language Adaptation}
    \subsubsection{Explainability}
    \lipsum[26][1-3]

    \subsubsection{Bias}
    \section{NLP in Education}
    \section{NLP for Accessibility}

    \section{Literature}
    \begin{itemize}
        \item Bird, S., Klein, E., \& Loper, E. (2009). Natural Language Processing with Python. O'Reilly.
        \item Jurafsky, D., \& Martin, J. H. (2020). Speech and Language Processing (3rd ed.). Prentice Hall. \url{https://web.stanford.edu/~jurafsky/slp3}
        \item Kamath, U., Liu, J., \& Whitaker, J. (2019). Deep Learning for NLP and Speech Recognition: Practical NLP, Speech, and Deep Learning using Python-based Open Source Tools. Springer.
        \item Bocklisch, T., Faulker, J., Pawlowski, N., \& Nichol, A. (2017). Rasa: Open Source Language Understanding and Dialogue Management. NIPS Workshop on Conversational AI.
        \item Chen, L., Chen, P., \& Lin, Z. (2020): Artificial Intelligence in Education: A Review. IEEE Access 8 (2020), 75264--75278.
        \item Heffernan, N. (2014): The ASSISTments Ecosystem: Building a Platform that Brings Scientists and Teachers Together for Minimally Invasive Research on Human Learning and Teaching. International Journal of Artificial Intelligence in Education 24 (12 2014).
        \item Libbrecht, P., Declerck, T., Schlippe, T., Mandl, T., \& Schiffner, D. (2020): NLP for Student and Teacher: Concept for an AI based Information Literacy Tutoring System. In The 29th ACM International Conference on Information and Knowledge Management (CIKM2020). Galway, Ireland.
        \item Schlippe, T., \& Sawatzki, J. (2021): AI-based Multilingual Interactive Exam Preparation. In The Learning Ideas Conference 2021 (14th annual conference). ALICE - Special Conference Track on Adaptive Learning via Interactive, Collaborative and Emotional Approaches. New York, USA.
        \item Schlippe, T., \& Sawatzki, J. (2021): Cross-Lingual Automatic Short Answer Grading. In Proceedings of The 2nd International Conference on Artificial Intelligence in Education Technology (AIET 2021). Wuhan, China.
        \item Al-Thanyyan, S. S., \& Azmi, A. (2021): Automated Text Simplification: A Survey. ACM Computing Surveys, Vol. 54, Issue 2, pp. 1--36.
        \item Klaper, D., Ebling, S., \& Volk, M. (2013): Building a German/Simple German Parallel Corpus for Automatic Text Simplification. The Second Workshop on Predicting and Improving Text Readibility for Target Reader Populations.
        \item Schlippe, T., Alessai, S., El-Taweel, G., Wölfel, M., \& Zaghouani, W. (2020): Visualizing Voice Characteristics with Type Design in Closed Captions for Arabic. Cyberworlds. Caen, France.
    \end{itemize}

